\section{Introduction}
\label{sec:intro}

A host intrusion detection rule is a small program with a large view. It
sees every process a machine starts, every file it opens and every
connection it makes, and it decides which of them deserve an alert. In the
monitors that run on most Linux fleets today the rule is also
\emph{trusted}: Falco evaluates its rules and loads its plugins inside one
privileged agent~\cite{falco}, Tetragon compiles its policies into the
kernel and a privileged agent~\cite{tetragon}, and a mistake in a rule,
or a rule from a source that should not have been trusted, runs with the
agent's view and the agent's power. Rules come from many places---vendor
packs, community repositories, a security team's own detectors, and
increasingly a model trained by someone else---and they change far more
often than the agent does.

This paper asks a narrow question about that arrangement. Suppose the
rules were untrusted and ran inside a WebAssembly sandbox, each in its own
instance, with no access to the host beyond the event it is handed and
with a memory cap, an instruction budget and a wall-clock deadline on
every call. \emph{What does that cost, and what does it buy?} The
question is narrow on purpose. WebAssembly's isolation properties are
established~\cite{wasm-spec,rlbox,wasmboxc}; eBPF-based collection is
established~\cite{falco,tetragon,tracee}; what is missing is a measurement
of the boundary between them in this particular use, where the sandboxed
code is fed by the kernel at tens of thousands of events per second and
the alternative it is compared against is not ``no sandbox'' but two
production systems.

We built \sys to take that measurement. Six fixed eBPF programs produce a
fixed \SI{592}{B} record per event; a Go daemon queues the records and
hands each one to every rule through a five-function ABI; a rule is a Wasm
module that returns a severity. Two rules ship with it: a pattern rule of
the kind one writes in Falco's filter language, and a syscall-sequence
detector that combines a STIDE $n$-gram database~\cite{stide} with a
small MLP whose weights are embedded at build time. The design choices
that affect the measurement are few: one long-lived instance per rule,
limits attached to the instance's store, a trapped rule dropped and
re-instantiated, and a choice between a JSON encoding of the event and the
kernel record passed verbatim.

The evaluation makes three claims, each with the experiment that supports
it.

\begin{enumerate}
\item \textbf{Cost (E2--E5).} A sandboxed rule call costs
  \eTwoEngineOverheadUs{}~\si{\micro\second} empty and
  \eTwoMlRawUs{}~\si{\micro\second} for the real detector on a process
  event; end to end, the pipeline sustains \eThreeAchievedBinFiftyK{}
  syscall events per second with the binary encoding
  (\eThreeAchievedJsonFiftyK{} with JSON) at a median latency of
  \eThreePFiftyBinTenK{}~\si{\micro\second} at \num{10000}~events/s. Under
  nginx, redis and PostgreSQL, \sys costs the application
  \eFourDeltaNginxWs{}, \eFourDeltaRedisWs{} and \eFourDeltaPgWs{}\% of
  throughput, in the range of Falco's \eFourDeltaNginxFalco{},
  \eFourDeltaRedisFalco{} and \eFourDeltaPgFalco{}\%, with a smaller
  resident footprint and, under PostgreSQL, a larger CPU one; the three
  systems detect the same three actions in every trial with alert
  latencies within a millisecond of each other. The CPU gap has a named
  cause---every event reaches every rule, and the ML rule recomputes its
  features per event---and is a lever on the rule side, not a sandbox
  tax.
\item \textbf{Isolation (E6).} \isoModules{} modules written to read
  files, open sockets, leak the environment, flood output, exhaust
  memory, spin forever, overflow the stack or exit the process were all
  contained under the default limits in \isoContained{} of \isoTotal{}
  attempts, with the neighbouring rule unaffected and the next call
  succeeding. We report the one capability a rule does have (the clock)
  and the one case the runtime classifies less precisely than we would
  like.
\item \textbf{Equivalence (E1--E2).} The module the daemon loads produces,
  to within \eTwoMaxDiffNumpyWasm{}, the same score as the reference
  implementation that was trained, on all \eTwoTraces{} ADFA-LD traces.
  Because of that, the accuracy we report is the module's own: every test
  score in a \eOneFolds-fold, \eOneSeeds-seed evaluation is produced
  inside the sandbox. That accuracy is an AUPRC of \eOneAuprcMlp{},
  above STIDE and an Isolation Forest and \emph{below} a random forest
  (\eOneAuprcRf{}) and an LSTM (\eOneAuprcLstm{}) on the same data. We
  say so because the point of E1 is the protocol, and because nothing in
  the sandbox prefers the weaker model.
\end{enumerate}

We do not claim that \sys is the earliest system to combine eBPF and
WebAssembly, nor that its detector is state of the art, nor that its
event path is copy-free. The nearest prior work combines the two
technologies in other directions (\S\ref{sec:related}); the detector is
deliberately ordinary; and the binary encoding removes a parse, not a
copy. What the paper contributes is the measurement, taken on the deployed
artifact, with every number regenerated from the result files we release
(\S\ref{sec:openscience}).
